\subsection{p-ADIC UPPER BOUNDS OF THE SERIES exp, log AND H}

In this no. we assume that \(p \neq 0\) .

Lemma 1. Let n be an integer \(\geqslant0\) and let \(n=n_{0}+n_{1}p+\cdots+n_{k}p^{k}\) , with \(0\leqslant n_{i}\leqslant p-1\) , be the p-adic expansion of n.~Let \(\mathrm{S}(n)=n_{0}+n_{1}+\cdots+n_{k}\) . Then

\[
v _ {p} (n!) = \frac {n - \mathrm{S} (n)}{p - 1}.\tag{3}
\]

\(v_{p}(n!) = \sum_{i=1}^{n} v_{p}(i)\) and the number of integers \(i\) between 1 and \(n\) for which \(v_{p}(i) \geqslant j\) is equal to the integral part \([n/p^{j}]\) of \(n/p^{j}\) . Then

\[
v _ {p} (n!) = \sum_ {j \geqslant 0} j ([ n / p ^ {j} ] - [ n / p ^ {j + 1} ]) = \sum_ {j \geqslant 1} [ n / p ^ {j} ].
\]

As \([n / p^j] = \sum_{i\geqslant j}n_i p^{i - j}\) , the lemma follows.

Lemma 2. \(v(n) \leqslant v(n!) \leqslant (n - 1)\theta\) and \(v(n) \leqslant (\log n) / (\log p)\) for every integer \(n \geqslant 1\) .

\(v(n!) = v_{p}(n!) = (n - S(n))\theta \leqslant (n - 1)\theta\) by Lemma 1.

On the other hand, \(n \geqslant p^{v(n)}\) , whence \(v(n) \leqslant (\log n) / (\log p)\) .

Let \(\mathbf{I} = \{\mathbf{U},\mathbf{V}\}\) be a set of two elements and let

\[
\mathrm{H} = \sum_ {r, s \geqslant 0} \mathrm{H} _ {r, s} (\mathrm{U}, \mathrm{V}) \in \hat {\mathrm{L}} _ {\mathbf {Q}} (\mathrm{I})
\]

be the Hausdorff series (§ 6, no. 4, Definition 1). Let \(\mathbf{Z}_{(p)}\) be the local ring of \(\mathbf{Z}\) relative to the prime ideal \((p)\) and \((e_b)_{b \in \mathbf{B}}\) a basis of \(\mathrm{L}_{\mathbf{Z}_{(p)}}(\mathrm{I})\) over \(\mathbf{Z}\) (§ 2, no. 11, Theorem 1). It is also a basis of \(\mathrm{L}_{\mathbf{Q}}(\mathrm{I})\) over \(\mathbf{Q}\) .

PROPOSITION 1. Let \(r\) and \(s\) be two integers \(\geqslant 0\) . If \(H_{r,s} = \sum_{b \in B} \lambda_b e_b\) , where \(\lambda_b \in \mathbf{Q}\) , is the decomposition of \(H_{r,s}\) with respect to the basis \((e_b)_{b \in B}\) , then

\[
v _ {p} (\lambda_ {b}) \geqslant - (r + s - 1) \theta \quad for all b \in \mathbf {B}.\tag{4}
\]

The ring \(\mathrm{A}_{\mathbf{Z}_{(p)}}(\mathrm{I})\) is identified with the sub- \(\mathbf{Z}_{(p)}\) -module of \(\mathrm{A}_{\mathbf{Q}}(\mathrm{I})\) generated by the words \(w \in \mathrm{Mo}(\mathrm{I})\) . As \(\mathrm{L}_{\mathbf{Z}_{(p)}}(\mathrm{I})\) is a direct factor of \(\mathrm{A}_{\mathbf{Z}_{(p)}}(\mathrm{I})\) ,

\[
\mathrm{L} _ {\mathbf {Z} _ {(p)}} (\mathrm{I}) = \mathrm{A} _ {\mathbf {Z} _ {(p)}} (\mathrm{I}) \cap \mathrm{L} _ {\mathbf {Q}} (\mathrm{I}).\tag{5}
\]

Let \(f\) be the integer such that \(f \leqslant (r + s - 1)\theta < f + 1\) . Relation (4) is equivalent to \(v_{p}(\lambda_{b}) \geqslant -f\) for all \(b \in B\) , that is \(H_{r,s} \in p^{-f}L_{\mathbf{Z}_{(p)}}(I)\) . But this is equivalent also by (5) to \(H_{r,s} \in p^{-f}A_{\mathbf{Z}_{(p)}}(I)\) .

By formula (11) of §6, no. 4, it suffices to show that, for every integer \(m \geqslant 1\) and all integers \(r_1, \ldots, r_m, s_1, \ldots, s_m\) such that

\[
r _ {1} + \dots + r _ {m} = r, \qquad s _ {1} + \dots + s _ {m} = s,\tag{6}
\]

\[
r _ {i} + s _ {i} \geqslant 1 \quad \text { for } 1 \leqslant i \leqslant m,
\]

we have

\[
v _ {p} (m. r! \dots r _ {m}! s _ {1}! \dots s _ {m}!) \leqslant f.\tag{7}
\]

But by Lemma 2, \(v_{p}(r_{i}!s_{i}!) \leqslant (r_{i} + s_{i} - 1)\theta\) and \(v_{p}(m) \leqslant v_{p}(m!) \leqslant (m - 1)\theta\) ; the left hand side of (7) is therefore bounded above by

\[
\theta (m - 1 + \sum_ {i = 1} ^ {m} \left(r _ {i} + s _ {i} - 1\right)) = \theta (r + s - 1);
\]

as it is an integer, it is \(\leqslant f\) , which completes the proof.

